% ======================================================================
%  sec24.tex —— 第 24 节 紧合态射
%  底本：原书 pp.114–118（印刷页 = PDF 页）
%  原书标题：24. Proper morphisms
% ======================================================================

\section{紧合态射}\label{sec:24}

% ---------------------------------------------------------------- p.114
\subsection{闭态射}\label{sec:24.1}

\begin{definition}\label{def:24.1.1}
称拓扑空间的连续映射 $f\colon X \to Y$ \emph{闭}，若它把 $X$ 的闭子集都送到 $Y$ 的闭子集。
\end{definition}

\begin{example}\label{ex:24.1.2}
概形的闭浸入是闭态射。
\end{example}

\begin{lemma}\label{lem:24.1.3}
设 $f\colon X \to Y$ 为概形的拟紧态射。则下列条件等价：
\begin{enumerate}[label=(\roman*),leftmargin=2em,itemsep=0.25ex]
  \item 子集 $f(X)$ 在 $Y$ 中闭；
  \item 子集 $f(X)$ 对特殊化稳定。
\end{enumerate}
\end{lemma}

\begin{proof}
(i)$\Rightarrow$(ii) 显然。

(ii)$\Rightarrow$(i)：注意问题在 $Y$ 上局部，故可设 $Y$ 仿射。因 $f$ 拟紧，$X$ 拟紧。取 $X$ 的有限仿射开子概形覆盖 $X = \bigcup_{i \in I} U_i$。考虑不交并 $X' := \bigsqcup_{i \in I} U_i$ 与合成 $f'\colon X' \to X \to Y$。注意 $f(X') = f(X)$。故可把 $X$ 换成 $X'$，假设 $X$ 仿射。

设 $\phi\colon A \to B$ 为对应于 $f\colon X \to Y$ 的同态。可把 $A$ 换成 $A/\ker(\phi)$，假设 $\phi$ 单射。于是 $f(X)$ 包含对应于 $\Spec(A)$ 的极小素点的全部点。由 (ii) 得 $f(X) = Y$。特别地 (i) 成立。
\end{proof}

\begin{yizhu}
「对特殊化稳定」是说：若 $y \in f(X)$ 且 $y' \in \overline{\{y\}}$，则 $y' \in f(X)$。Stacks Project 把同一断言写作「$f(X)$ 闭」与「$f(X)$ 对特殊化稳定」，并把拟紧性用在一处：把 $X$ 换成有限个仿射开片的不交并，使像集不变，从而化到 $X$ 仿射。
[来源：Stacks Project, Tag 05JL（Lemma 26.19.7，拟紧态射的像），\url{https://stacks.math.columbia.edu/tag/05JL}（2026-11-17 访问）。]
\end{yizhu}

\begin{corollary}\label{cor:24.1.4}
设 $f\colon X \to Y$ 为概形的拟紧态射。则下列条件等价：
\begin{enumerate}[label=(\roman*),leftmargin=2em,itemsep=0.25ex]
  \item $f$ 闭；
  \item 特殊化沿 $f$ 提升。换言之，对 $y, y' \in Y$ 与 $x \in X$，若 $y' \in \overline{\{y\}}$ 且 $f(x) = y$，则存在 $x' \in \overline{\{x\}}$，使 $f(x') = y'$。
\end{enumerate}
\end{corollary}

\begin{proof}
由引理 \ref{lem:24.1.3} 与「$X$ 的任意闭子集都是某个闭子概形的底集」即得。
\end{proof}

\begin{yizhu}
这条引理是紧合性的拓扑骨架：闭性被翻译成「特殊化可沿 $f$ 提升」。Stacks Project 的对应条目正是「拟紧时 $f$ 闭，当且仅当特殊化沿 $f$ 提升」，证明把 $X$ 的闭子集看作某个闭子概形的底集，再对限制态射用前一条引理。
[来源：Stacks Project, Tag 01K9（Lemma 26.19.8，闭性与特殊化的提升），\url{https://stacks.math.columbia.edu/tag/01K9}（2026-11-17 访问）。]
\end{yizhu}

下列结果的证明留给读者。

\begin{lemma}\label{lem:24.1.5}
考虑概形的闭态射构成的类。
\begin{enumerate}[label=(\roman*),leftmargin=2em,itemsep=0.25ex]
  \item 设 $X \xrightarrow{f} Y \xrightarrow{g} Z$ 为态射链，$f$ 与 $g$ 都闭。则 $g \circ f$ 亦然。
  \item 「闭」在目标上是局部的。
\end{enumerate}
\end{lemma}

\subsection{泛闭态射}\label{sec:24.2}

\begin{definition}\label{def:24.2.1}
称概形态射 $f\colon X \to Y$ \emph{泛闭}，若对任意 $Y$ 上的概形 $Y'$，基变换 $f'\colon X \times_Y Y' \to Y'$ 都闭。
\end{definition}

\begin{yizhu}
「泛」字的含义是对一切基变换都要求闭性，故泛闭是基变换稳定的性质，而单纯的闭不是。这解释了紧合性为何用泛闭而非闭：$\A^1_k \to \Spec(k)$ 闭，作一次基变换后却不闭（例 \ref{ex:24.2.2}），故它不该被认作紧合。
[来源：Stacks Project, Tag 01KB（Definition 26.20.1，泛闭态射），\url{https://stacks.math.columbia.edu/tag/01KB}（2026-11-17 访问）。]
\end{yizhu}

\begin{example}\label{ex:24.2.2}
设 $k$ 为域。态射 $f\colon \A^1 \to \Spec(k)$ 闭，但其基变换 $\A^2 \to \A^1$ 不闭。事实上，闭子集 $Z((xy-1)) \subseteq \A^2$ 的像是 $\A^1 \setminus 0 \subseteq \A^1$。
\end{example}

% ---------------------------------------------------------------- p.115
\begin{lemma}\label{lem:24.2.3}
概形的泛闭态射 $f\colon X \to Y$ 拟紧。
\end{lemma}

\begin{proof}
见 [Sta25, Lemma 04XU]。
\end{proof}

\begin{yizhu}
该引理出自 Poonen，证明避开 Zorn 引理与诺特假设。若 $f$ 不拟紧，取仿射开子集 $V \subseteq Y$ 使 $f^{-1}(V)$ 不拟紧，再用无限多个仿射开片构造一个基变换，使像集不是闭集。原文只引 Tag，未给证明。
[来源：Stacks Project, Tag 04XU（Lemma 29.42.8，泛闭态射拟紧），\url{https://stacks.math.columbia.edu/tag/04XU}（2026-11-17 访问）。]
\end{yizhu}

下列结果的证明留给读者。

\begin{lemma}\label{lem:24.2.4}
考虑概形的泛闭态射构成的类。
\begin{enumerate}[label=(\roman*),leftmargin=2em,itemsep=0.25ex]
  \item 设 $X \xrightarrow{f} Y \xrightarrow{g} Z$ 为态射链，$f$ 与 $g$ 都泛闭。则 $g \circ f$ 亦然。
  \item 设 $X \xrightarrow{f} Y \xrightarrow{g} Z$ 为态射链，$g \circ f$ 泛闭且 $g$ 分离。则 $f$ 泛闭。
  \item 泛闭态射对基变换稳定。
  \item 「泛闭」在目标上是局部的。
\end{enumerate}
\end{lemma}

\subsection{赋值判别法}\label{sec:24.3}

回忆下述定义。

\begin{definition}\label{def:24.3.1}
\emph{赋值环}是指满足下述条件的整环 $A$：对其分式域 $K := \Frac(A)$ 的任意非零元 $x$，$x$ 与 $x^{-1}$ 中至少有一个属于 $A$。
\end{definition}

\begin{yizhu}
nLab 在同一页列出等价的三种刻画：$x$ 或 $x^{-1}$ 属于 $A$；理想按包含构成全序；主理想按包含构成全序。由全序性立得不可逆元之集对加法封闭，这正是习题 \ref{exc:24.3.2} 的内容，故赋值环自动局部。
[来源：nLab "valuation ring"，\url{https://ncatlab.org/nlab/show/valuation+ring}（Last revised 2025-07-29；2026-11-17 访问）。]
\end{yizhu}

\begin{exercise}\label{exc:24.3.2}
赋值环 $A$ 是局部环。
\end{exercise}

再回忆交换代数中的下述事实。

\begin{lemma}\label{lem:24.3.3}
设 $K$ 为域，$\mathcal{C}$ 为 $K$ 的局部子环构成的集合。对 $A, B \in \mathcal{C}$，称 $B \le A$，若 $B \subseteq A$ 且 $B \to A$ 是局部同态。则对 $A \in \mathcal{C}$，下列条件等价：
\begin{itemize}
  \item $A$ 是偏序集 $(\mathcal{C}, \le)$ 的极大元；
  \item $A$ 是赋值环，且典范映射 $\Frac(A) \to K$ 是同构。
\end{itemize}
\end{lemma}

\begin{definition}\label{def:24.3.4}
设 $f\colon X \to Y$ 为概形态射。称 $f$ 满足赋值判别法的\emph{存在部分}，若对任意赋值环 $A$（记 $K = \Frac(A)$）与任意实线交换方块
\[
\begin{tikzcd}
\Spec(K) \arrow[r] \arrow[d] & X \arrow[d, "f"] \\
\Spec(A) \arrow[r] \arrow[ur, dashed] & Y,
\end{tikzcd}
\tag{24.1}\label{eq:24.1}
\]
都存在虚线态射使两个三角形交换。

称 $f$ 满足赋值判别法的\emph{唯一部分}，若在上述情形下，使两个三角形交换的虚线态射至多一个。

称 $f$ 满足\emph{赋值判别法}，若它同时满足存在部分与唯一部分。
\end{definition}

\begin{proposition}\label{prop:24.3.5}
设 $f\colon X \to Y$ 为概形的拟紧态射。则 $f$ 泛闭，当且仅当它满足赋值判别法的存在部分。
\end{proposition}

% ---------------------------------------------------------------- p.116
\begin{proof}
先证「仅当」。设 $f$ 泛闭。对给定的实线交换图 \eqref{eq:24.1}，考虑 $X_A := X \times_Y \Spec(A)$。容易看出：知道使 \eqref{eq:24.1} 交换的虚线态射，等价于知道使下图交换的虚线态射：
\[
\begin{tikzcd}
\Spec(K) \arrow[r] \arrow[d] & X_A \arrow[d, "f_A"] \\
\Spec(A) \arrow[r, equal] & \Spec(A).
\end{tikzcd}
\]
因此可把 $f$ 换成 $f_A$，从而假设 $\Spec(A) \to Y$ 是恒等态射。

设 $x \in X$ 为 $\Spec(K) \to X$ 的像。注意 $f(x)$ 是 $\Spec(A)$ 的一般点。由推论 \ref{cor:24.1.4}，存在 $x' \in \overline{\{x\}}$，使 $f(x')$ 是 $\Spec(A)$ 的闭点。只需证存在使下图交换的虚线态射：
\[
\begin{tikzcd}
\Spec(K) \arrow[r] \arrow[d] & \Spec(\cO_{X,x'}) \arrow[d] \\
\Spec(A) \arrow[r, equal] & \Spec(A).
\end{tikzcd}
\]
注意右侧纵向态射把闭点送到闭点，故下图右侧的环同态是局部同态：
\[
\begin{tikzcd}
K & \cO_{X,x'} \arrow[l] \\
A \arrow[u] & A \arrow[u] \arrow[l, "\cong"']
\end{tikzcd}
\]
令 $B := \im(\cO_{X,x'} \to K)$。容易看出 $B$ 是局部环，且合成 $A \to \cO_{X,x'} \to B$ 是局部同态。由引理 \ref{lem:24.3.3}，同态 $A \to B$ 是同构。于是合成 $\cO_{X,x'} \to B \simeq A$ 给出所需的虚线态射。

再证「当」。设 $f$ 满足赋值判别法的存在部分。形式上可知 $f$ 的任意基变换也满足同一性质。故只需证 $f$ 闭。由推论 \ref{cor:24.1.4}，只需证特殊化沿 $f$ 提升。取 $y, y' \in Y$ 与 $x \in X$，使 $y' \in \overline{\{y\}}$ 且 $f(x) = y$。

注意有典范交换图
\[
\begin{tikzcd}
\Spec(\kappa_x) \arrow[r] \arrow[d] & X \arrow[d, "f"] \\
\Spec(\cO_{Y,y'}) \arrow[r] & Y,
\end{tikzcd}
\]
其中左侧纵向态射由同态
\[
\cO_{Y,y'} \lto \cO_{Y,y} \lto \kappa_y \lto \kappa_x \tag{24.2}\label{eq:24.2}
\]
给出。

% ---------------------------------------------------------------- p.117
记 $K := \kappa_x$，并设 $B \subseteq K$ 为 \eqref{eq:24.2} 的像。注意 $B$ 是局部环。由引理 \ref{lem:24.3.3}，可取赋值子环 $A \subseteq K$，使 $\Frac(A) \xrightarrow{\simeq} K$，且 $B \subseteq A$、$B \to A$ 是局部同态。特别地，\eqref{eq:24.2} 等于合成
\[
\cO_{Y,y'} \lto A \lto K,
\]
其中第一个态射是局部同态。这给出交换图
\[
\begin{tikzcd}
\Spec(K) \arrow[rr, dashed, "g"] \arrow[d] & & X \arrow[d, "f"] \\
\Spec(A) \arrow[r] & \Spec(\cO_{Y,y'}) \arrow[r] & Y.
\end{tikzcd}
\]
按构造，底行的合成把闭点送到 $y'$、把一般点送到 $y$。由假设，存在使图交换的虚线箭头 $g$。于是 $g$ 把闭点送到某个点 $x' \in X$，使 $x' \in \overline{\{x\}}$ 且 $f(x') = y'$。故特殊化沿 $f$ 提升，如所欲证。
\end{proof}

\begin{yizhu}
Stacks Project 把这一命题题为「泛闭性的赋值判别法」，并指出它是两条引理的直接推论。拟紧假设只用在「当」的方向：从存在部分回到泛闭，须借推论 \ref{cor:24.1.4} 把闭性化为特殊化提升；反方向里泛闭本身已蕴含拟紧（引理 \ref{lem:24.2.3}）。
[来源：Stacks Project, Tag 01KF（Proposition 26.20.6，泛闭性的赋值判别法），\url{https://stacks.math.columbia.edu/tag/01KF}（2026-11-17 访问）。]
\end{yizhu}

\begin{proposition}\label{prop:24.3.6}
设 $f\colon X \to Y$ 为概形的拟分离态射。则 $f$ 分离，当且仅当它满足赋值判别法的唯一部分。
\end{proposition}

\begin{proof}
因 $f$ 拟分离，$\Delta_f\colon X \to X \times_Y X$ 是拟紧浸入。于是有下述蕴含链：
\[
\begin{aligned}
& f \text{ 分离；}\\
\Leftrightarrow\ & \Delta_f \text{ 是闭浸入；}\\
\Leftrightarrow\ & \Delta_f \text{ 的任意基变换是闭浸入；}\\
\Leftrightarrow\ & \Delta_f \text{ 泛闭；}\\
\Leftrightarrow\ & X \to X \times_Y X \text{ 满足赋值判别法的存在部分；}\\
\Leftrightarrow\ & X \to Y \text{ 满足赋值判别法的唯一部分。}
\end{aligned}
\]
\end{proof}

\begin{yizhu}
Stacks Project 的分离性赋值判别法只陈述链条的最后一环：若 $f$ 拟分离且满足唯一部分，则 $f$ 分离。链条中其余各步都靠「$\Delta_f$ 是浸入」与泛闭性的赋值判别法。原书把这几步并成一行，便于整体记忆。
[来源：Stacks Project, Tag 01L0（Lemma 26.22.2，分离性的赋值判别法），\url{https://stacks.math.columbia.edu/tag/01L0}（2026-11-17 访问）。]
\end{yizhu}

利用诺特逼近，可以证明命题 \ref{prop:24.3.5} 与命题 \ref{prop:24.3.6} 的如下变体。

\begin{proposition}\label{prop:24.3.7}
设 $f\colon X \to Y$ 为概形的有限型态射，且 $Y$ 局部诺特。则 $f$ 泛闭，当且仅当它对所有离散赋值环满足赋值判别法的存在部分。
\end{proposition}

\begin{proof}
见 [Sta25, Lemma 05JY]。
\end{proof}

\begin{proposition}\label{prop:24.3.8}
设 $f\colon X \to Y$ 为概形的局部有限型态射，且 $Y$ 局部诺特。则 $f$ 分离，当且仅当它对所有离散赋值环满足赋值判别法的唯一部分。
\end{proposition}

\begin{proof}
见 [Sta25, Lemma 0207]。
\end{proof}

\begin{remark}\label{rem:24.3.9}
赋值环诺特，当且仅当它是离散赋值环。
\end{remark}

\subsection{紧合态射}\label{sec:24.4}

\begin{definition}\label{def:24.4.1}
称概形态射 $f\colon X \to Y$ \emph{紧合（proper）}，若它分离、有限型且泛闭。
\end{definition}

\begin{yizhu}
nLab 把紧合与拓扑紧性直接类比：紧合概形之于概形，如同紧空间之于拓扑空间；「拟紧」这一命名正是为避开与「紧」混淆。该条目还提到紧合性有 Chevalley 形式的赋值判别法，与命题 \ref{prop:24.3.5} 相呼应。
[来源：nLab "proper morphism"，\url{https://ncatlab.org/nlab/show/proper+morphism}（Last revised 2024-08-24；2026-11-17 访问）。]
\end{yizhu}

% ---------------------------------------------------------------- p.118
\begin{remark}\label{rem:24.4.2}
由引理 \ref{lem:24.2.3}，可把第二个条件换成「局部有限型」。
\end{remark}

\begin{example}\label{ex:24.4.3}
概形的任意闭浸入都紧合。
\end{example}

\begin{exercise}\label{exc:24.4.4}
证明 $\PP^n_{\Z} \to \Spec(\Z)$ 紧合。
\end{exercise}

下列结果由命题 \ref{prop:24.3.5} 与命题 \ref{prop:24.3.6} 合并得到。

\begin{corollary}\label{cor:24.4.5}
设 $f\colon X \to Y$ 为概形的拟分离有限型态射。则 $f$ 紧合，当且仅当它满足赋值判别法。
\end{corollary}

诺特情形下还有如下变体。

\begin{corollary}\label{cor:24.4.6}
设 $f\colon X \to Y$ 为概形的有限型态射，且 $Y$ 局部诺特。则 $f$ 紧合，当且仅当它对所有离散赋值环满足赋值判别法。
\end{corollary}

下列结果的证明留给读者。

\begin{lemma}\label{lem:24.4.7}
考虑概形的紧合态射构成的类。
\begin{enumerate}[label=(\roman*),leftmargin=2em,itemsep=0.25ex]
  \item 设 $X \xrightarrow{f} Y \xrightarrow{g} Z$ 为态射链，$f$ 与 $g$ 都紧合。则 $g \circ f$ 亦然。
  \item 设 $X \xrightarrow{f} Y \xrightarrow{g} Z$ 为态射链，$g \circ f$ 紧合且 $g$ 分离。则 $f$ 紧合。
  \item 紧合态射对基变换稳定。
  \item 「紧合」在目标上是局部的。
\end{enumerate}
\end{lemma}

下述定理将在《代数几何 II》中证明。

\begin{theorem}\label{thm:24.4.8}
设 $f\colon X \to Y$ 为有限表现的紧合态射。则 $f_*$ 保持凝聚模。
\end{theorem}

\begin{remark}\label{rem:24.4.9}
事实上，导出函子 $R^i f_*$ 也保持凝聚模。
\end{remark}

\begin{warning}\label{warn:24.4.10}
上述定理中「有限表现」这一条件必要。
\end{warning}

\begin{exercise}\label{exc:24.4.11}
设 $k$ 为域。证明：沿闭浸入 $\{0\} \into \A^{\infty}_k$ 的直接像函子不保持凝聚模。
\end{exercise}
